% piTantum -- Manuale generale (Parte I)
% Per dirigenti scolastici, coordinatori d'orario, docenti utenti.
% ASCII puro, accenti via macro \`e \'e \`a, virgolette ``...''.
% Manuale tecnico del progetto piTantum (Tempus Tantum;
% codename interno: timetable).
% ASCII puro, accenti via macro \`e \'e \`a, virgolette ``...''.
\documentclass[a4paper,11pt,oneside]{report}

\usepackage[T1]{fontenc}
\usepackage[utf8]{inputenc}
\usepackage[italian]{babel}
\usepackage{lmodern}
\usepackage{microtype}
\usepackage{geometry}
\geometry{margin=2.4cm,headheight=14pt}
\usepackage{xcolor}
\usepackage{graphicx}
\usepackage{tikz}
\usetikzlibrary{positioning,arrows.meta,shapes.geometric,fit,calc}
\usepackage{listings}
\usepackage{enumitem}
\usepackage{array}
\usepackage{longtable}
\usepackage{booktabs}
\usepackage{fancyhdr}
\usepackage{titlesec}
\usepackage[hidelinks,bookmarksopen=true]{hyperref}

\definecolor{cHard}{RGB}{220,38,38}
\definecolor{cSoft}{RGB}{217,119,6}
\definecolor{cPref}{RGB}{37,99,235}
\definecolor{cEnf}{RGB}{6,95,70}
\definecolor{cAllow}{RGB}{16,185,129}
\definecolor{cBox}{RGB}{229,231,235}
\definecolor{cCode}{RGB}{30,41,59}

\lstset{
  basicstyle=\ttfamily\footnotesize,
  keywordstyle=\color{cPref}\bfseries,
  commentstyle=\color{gray}\itshape,
  stringstyle=\color{cEnf},
  showstringspaces=false,
  breaklines=true,
  breakatwhitespace=true,
  columns=fullflexible,
  frame=single,
  rulecolor=\color{cBox},
  framesep=4pt,
  xleftmargin=8pt,
  xrightmargin=4pt,
  numbers=none,
}
\lstdefinestyle{json}{language=,morekeywords={true,false,null}}
\lstdefinestyle{py}{language=Python}
\lstdefinestyle{shell}{language=bash,morekeywords={apt,brew,dnf,pacman,pip,npm,git,curl}}

\titleformat{\chapter}[hang]{\Huge\bfseries}{\thechapter.}{12pt}{}
\titleformat{\section}{\Large\bfseries}{\thesection}{10pt}{}
\titleformat{\subsection}{\large\bfseries}{\thesubsection}{8pt}{}

\pagestyle{fancy}
\fancyhf{}
\fancyhead[L]{\textit{piTantum -- manuale tecnico}}
\fancyhead[R]{\thepage}
\renewcommand{\headrulewidth}{0.4pt}

\title{}
\author{}
\date{}



\begin{document}

% =========== Cover (Parte I) ===========
\begin{titlepage}
\thispagestyle{empty}
\centering
\vspace*{1.4cm}

{\Huge\bfseries $\pi$\,Tantum\par}
\vspace{0.2cm}
{\Large alias \emph{Tempus Tantum}\par}
\vspace{0.4cm}
{\Large\itshape sistema di generazione orario scolastico\par}

\vspace{1.0cm}
\rule{0.5\textwidth}{0.4pt}

\vspace{0.8cm}
\begin{minipage}{0.85\textwidth}
\centering
\itshape
``Omnia, Lucili, aliena sunt, tempus tantum nostrum est.''\\
\upshape\small Seneca, \emph{Epistulae morales ad Lucilium}, I, 1
\end{minipage}

\vspace{1.0cm}
\rule{0.5\textwidth}{0.4pt}

\vspace{1.4cm}
{\LARGE Manuale generale\par}
\vspace{0.4cm}
{\Large\itshape Parte I --- per coordinatori d'orario\par}

\vspace{2cm}

\begin{minipage}{0.8\textwidth}
\centering\small
Una guida discorsiva all'uso quotidiano del sistema:
preparazione dell'orario di inizio anno, gestione delle assenze,
modifica in corso d'anno, creazione di corsi di recupero.
Pensata per chi usa $\pi$Tantum come strumento di lavoro,
senza dover conoscerne i dettagli implementativi.
\end{minipage}

\vfill

{\large Giovanni Borghi\par}
\vspace{0.2cm}
{\large \today\par}
\end{titlepage}

\tableofcontents

% =========== Capitoli generali ===========
\chapter{Introduzione}

\section*{Come usare questo manuale}

La documentazione di $\pi$Tantum \`e suddivisa in due parti, per
adattarsi a pubblici diversi:

\begin{itemize}
\item Il \textbf{Manuale generale} (Parte I, file
  \texttt{manual\_generale.pdf}) \`e pensato per chi usa
  $\pi$Tantum come strumento di lavoro: dirigenti scolastici,
  coordinatori d'orario, docenti che devono modificare un
  vincolo o capire perch\'e un orario \`e venuto in un certo
  modo. Tono discorsivo, casi d'uso concreti, niente jargon.
\item Le \textbf{Appendici tecniche} (Parte II, file
  \texttt{appendici\_tecniche.pdf}) sono per sviluppatori,
  amministratori IT e curiosi: architettura interna, modello
  dati, API REST, grammatica formale del DSL, dettaglio delle
  strategie di ottimizzazione e della diagnostica statistica.
\item Il \textbf{Manuale completo} (file
  \texttt{manual\_completo.pdf}) raccoglie entrambe le parti
  in un unico volume, utile per chi vuole stamparlo o averlo
  tutto a portata di mano.
\end{itemize}

Tutti e tre i PDF sono ricostruibili dai sorgenti
\texttt{docs/*.tex} con lo script \texttt{docs/build\_manual.sh}
(o \texttt{.bat} su Windows).

\section*{Cos'\`e $\pi$Tantum}

\emph{$\pi$Tantum} (alias \emph{Tempus Tantum}, codename interno:
\texttt{timetable}) \`e un sistema di generazione e gestione
dell'orario scolastico per un Liceo italiano. Il nome richiama il
celebre incipit dell'epistola I di Seneca a Lucilio:
``Omnia, Lucili, aliena sunt, tempus tantum nostrum est'' (Seneca,
\emph{Epistulae morales ad Lucilium}, I, 1) -- ``tutto ci appartiene
di altri, soltanto il tempo \`e nostro''. Il gioco grafico fra la
lettera greca $\pi$ e le due ``T'' di \emph{Tempus Tantum} riassume
l'idea: organizzare il tempo della scuola, dato che del tempo non
si dispone se non gestendolo.

\`E composto di tre layer:

\begin{enumerate}
  \item \textbf{Solver} (\texttt{experiments/}): pipeline CP-SAT
        (Google OR-Tools) con decomposizione spettrale per istanze
        grandi e cascata metaeuristica (LNS, SA, TS, ILS).
  \item \textbf{Web UI} (\texttt{webui/}): backend FastAPI su
        SQLite + frontend SvelteKit. \`E il punto di accesso quotidiano:
        gestisce CRUD, ottimizzazione, visualizzazione, drag-and-drop,
        assenze e supplenze.
  \item \textbf{Legacy} (\texttt{schedule/}): script monolitici da cui
        il solver \`e nato; mantenuti per riferimento storico.
\end{enumerate}

Il manuale copre tutti e tre i layer ma si concentra sulla web UI,
che ha assorbito quasi ogni funzionalit\`a operativa.

\section{Avvio rapido}

\paragraph{Windows.} Doppio click su \texttt{webui/start.bat}.
Apre due finestre cmd: backend \emph{uvicorn} su porta 8000, frontend
\emph{vite} su porta 5173.

\paragraph{Linux / macOS.}
\begin{lstlisting}[style=shell]
cd timetable/webui
./start.sh           # background; URLs sulla console
./stop.sh            # ferma backend + frontend
\end{lstlisting}

Nel browser, aprire \texttt{http://127.0.0.1:5173} e dalla
\textbf{Dashboard} importare un profilo g\`ia calcolato
(\texttt{small} per cominciare).

\section{Stato attuale del progetto}

Le funzionalit\`a principali, al momento di questa redazione:

\begin{itemize}
  \item 16 sezioni (tab) della UI, $\sim$118 endpoint REST.
  \item $\sim$30 tabelle SQLAlchemy con migrazioni idempotenti.
  \item 5 stati di vincoli (HARD/SOFT/PREFERITO/ENFORCED/ALLOWED).
  \item Vincoli logici DNF con DSL esteso (predicate atoms
        \texttt{aula:}, \texttt{materia:}, \texttt{classe:},
        \texttt{gruppo:}, \texttt{docente:}).
  \item Conflict detection automatico nel tab \emph{Vincoli}.
  \item Drag-and-drop con preview live (rosso/giallo/verde) e
        \emph{room-follows-lesson}.
  \item Gestione assenze settimanali con assegnazione supplenti
        drag-drop.
  \item Tab \emph{Monitor} con espansione lezione-per-lezione e
        slot picker matriciale.
\end{itemize}

% =========== Cap 2: Installazione ===========

\chapter{Installazione}
\label{ch:installazione}

Guida cross-platform per installare ed eseguire piTantum su Windows,
Linux e macOS. La versione \emph{di riferimento} di questa guida \`e
\texttt{docs/installation.md} nel repo (Markdown, sempre aggiornato);
qui se ne riproduce il contenuto in forma compatta per il manuale
PDF. In caso di divergenza, fa fede il file Markdown.

\paragraph{Architettura del lancio.} Il progetto \`e composto da due
processi:
\begin{itemize}
  \item \textbf{Backend}: \texttt{uvicorn} che serve l'app FastAPI in
        \texttt{webui/backend} sulla porta 8000.
  \item \textbf{Frontend}: server di sviluppo \texttt{vite} che serve
        l'app SvelteKit in \texttt{webui/frontend} sulla porta 5173.
\end{itemize}
Uno script unico (\texttt{start.bat} su Windows,
\texttt{start.sh} su Linux/macOS) li avvia entrambi. Al primo lancio
crea la \emph{venv} Python, installa i requirements, fa
\texttt{npm install}; ai lanci successivi parte in pochi secondi.

% ----------------------------------------------------------------
\section{Windows}

\paragraph{Prerequisiti.}
\begin{itemize}
  \item Python $\geq$ 3.11 -- installer ufficiale:
        \url{https://www.python.org/downloads/windows/}
  \item Node.js $\geq$ 20 LTS -- installer ufficiale:
        \url{https://nodejs.org/en/download} (scegliere ``LTS'').
  \item Git -- installer ufficiale:
        \url{https://git-scm.com/download/win}.
\end{itemize}

\paragraph{Note importanti.}
\begin{itemize}
  \item Sull'installer Python, \textbf{spuntare} ``Add python.exe to
        PATH''. Senza, lo script non trova il comando \texttt{python}.
  \item Riavviare il terminale dopo le installazioni: il PATH non
        si aggiorna nelle finestre g\`ia aperte.
\end{itemize}

\paragraph{Clone + lancio.}
\begin{lstlisting}[style=shell]
git clone https://github.com/gborghi/timetable.git
cd timetable
webui\start.bat
\end{lstlisting}
Lo script apre due finestre cmd (``piTantum -- backend'' e
``piTantum -- frontend''). Aprire il browser su
\url{http://localhost:5173}.

\paragraph{Stop.} Chiudere le due finestre cmd, oppure
\texttt{Ctrl+C} in ciascuna.

\paragraph{Troubleshooting Windows.}
\begin{itemize}
  \item ``\texttt{python} non riconosciuto'': l'installer Python non
        \`e stato eseguito con ``Add to PATH''. Reinstallare oppure
        aggiungere \texttt{C:\textbackslash Users\textbackslash
        \emph{tu}\textbackslash AppData\textbackslash Local\textbackslash
        Programs\textbackslash Python\textbackslash Python311} al PATH.
  \item Porta 8000 occupata: \texttt{netstat -ano | findstr :8000}
        + \texttt{taskkill /PID <pid> /F}.
  \item Antivirus blocca \texttt{start.bat}: aggiungere alle
        eccezioni di Windows Defender.
\end{itemize}

% ----------------------------------------------------------------
\section{Linux (Ubuntu, Debian, Fedora, Arch)}

\paragraph{Prerequisiti.} Python $\geq$ 3.11, Node $\geq$ 20 LTS,
Git, build-essential.

\paragraph{Ubuntu / Debian.}
\begin{lstlisting}[style=shell]
sudo apt update
sudo apt install -y python3 python3-venv python3-pip nodejs npm \
                    git build-essential
\end{lstlisting}

I repo \texttt{apt} di Ubuntu 22.04 / Debian 12 hanno Node \emph{12-18},
troppo vecchio per SvelteKit 2 (richiede 20+). In quel caso usare
\textbf{nvm}:
\begin{lstlisting}[style=shell]
curl -o- https://raw.githubusercontent.com/nvm-sh/nvm/v0.39.7/install.sh | bash
source ~/.bashrc
nvm install --lts && nvm use --lts
\end{lstlisting}

\paragraph{Fedora.}
\begin{lstlisting}[style=shell]
sudo dnf install -y python3 python3-virtualenv nodejs npm git \
                    gcc make
\end{lstlisting}

\paragraph{Arch / Manjaro.}
\begin{lstlisting}[style=shell]
sudo pacman -S python python-pip nodejs npm git base-devel
\end{lstlisting}

\paragraph{Clone + lancio.}
\begin{lstlisting}[style=shell]
git clone https://github.com/gborghi/timetable.git
cd timetable
chmod +x webui/start.sh webui/stop.sh
./webui/start.sh                # background; log in webui/logs/
./webui/start.sh --foreground   # log live in console (Ctrl+C ferma)
\end{lstlisting}

\paragraph{Stop.}
\begin{lstlisting}[style=shell]
./webui/stop.sh
\end{lstlisting}

\paragraph{Troubleshooting Linux.}
\begin{itemize}
  \item Porta in uso: \texttt{sudo lsof -i :8000} + \texttt{kill <pid>}
        (eventualmente \texttt{kill -9 <pid>}).
  \item \texttt{python3} mancante su distro minimal: usare
        l'eseguibile chiamato \texttt{python} (3.x).
  \item WSL2: \texttt{start.sh} funziona normalmente. Aprire il
        browser dal lato Windows su
        \url{http://localhost:5173}; WSL2 inoltra le porte
        automaticamente.
\end{itemize}

% ----------------------------------------------------------------
\section{macOS (Intel + Apple Silicon)}

\paragraph{Prerequisiti.}
\begin{itemize}
  \item Xcode Command Line Tools: \texttt{xcode-select --install}.
  \item Homebrew: vedere \url{https://brew.sh/}.
  \item Python $\geq$ 3.11, Node $\geq$ 20 LTS, Git: tutti via Homebrew.
\end{itemize}

\paragraph{Installazione.}
\begin{lstlisting}[style=shell]
# Homebrew (se non presente)
/bin/bash -c "$(curl -fsSL https://raw.githubusercontent.com/Homebrew/install/HEAD/install.sh)"

# Toolchain
brew install python@3.12 node git

# Verifica
python3 --version    # >= 3.11
node --version       # >= v20
\end{lstlisting}

\paragraph{Note Apple Silicon.}
\begin{itemize}
  \item \texttt{ortools} dalla v9.7 spedisce wheel ARM64 nativi:
        Rosetta non serve.
  \item Se Python \`e installato da python.org invece che Homebrew,
        scegliere l'installer ``macOS 64-bit universal2''.
\end{itemize}

\paragraph{Clone + lancio + stop.} Identico a Linux.

\paragraph{Troubleshooting macOS.}
\begin{itemize}
  \item Gatekeeper blocca \texttt{start.sh}:
        \texttt{xattr -d com.apple.quarantine webui/start.sh
        webui/stop.sh}.
  \item \texttt{ortools} senza wheel ARM (versioni vecchie):
        \texttt{pip install --upgrade ortools} dentro la venv del
        backend.
  \item \texttt{brew} non in PATH:
        \texttt{eval "\$(/opt/homebrew/bin/brew shellenv)"} in
        \texttt{\textasciitilde/.zprofile}.
\end{itemize}

% ----------------------------------------------------------------
\section{Verifica installazione}

Dopo il primo lancio andato a buon fine:

\begin{lstlisting}[style=shell]
python --version       # >= 3.11   (Linux/macOS: spesso python3)
node --version         # >= v20
npm --version          # >= 10

# Backend health
curl http://127.0.0.1:8000/api/health
# atteso: {"status":"ok","name":"pitantum","version":"0.1.0"}

# Async layer (Section 2.5 P2)
curl http://127.0.0.1:8000/api/health/async
# atteso: {"status":"ok","async":true,"dialect":"sqlite"}

# Frontend
curl -I http://127.0.0.1:5173
# atteso: HTTP/1.1 200 OK
\end{lstlisting}

% ----------------------------------------------------------------
\section{Aggiornamento}

Quando si fa \texttt{git pull} per portare a casa nuovi commit:

\begin{lstlisting}[style=shell]
cd /path/to/timetable
git pull origin main

# se sono cambiate dipendenze Python:
webui/backend/.venv/bin/pip install -r webui/backend/requirements.txt
# (Windows: webui\backend\.venv\Scripts\pip install ...)

# se sono cambiate dipendenze frontend:
cd webui/frontend && npm install && cd ../..

# se ci sono migration DB:
cd webui/backend && ./.venv/bin/alembic upgrade head && cd ../..
# (Windows: webui\backend\.venv\Scripts\alembic upgrade head)
\end{lstlisting}

In pratica \texttt{start.sh}/\texttt{start.bat} rilanciato dopo un
\texttt{git pull} gestisce automaticamente \texttt{pip install} e
\texttt{npm install} quando servono. Le migration DB sono coperte da
\texttt{\_apply\_lightweight\_migrations()} allo startup come
safety-net (vedi capitolo \emph{Modello dati}); per la via canonica
usare \texttt{alembic upgrade head}.

Dopo un aggiornamento backend, fare un \textbf{hard reload} del
browser (\texttt{Ctrl+Shift+R}) per pulire chunk JavaScript cached.

% ----------------------------------------------------------------
\section{Disinstallazione}

\begin{lstlisting}[style=shell]
cd /path/to/timetable
./webui/stop.sh                 # Linux/macOS
# (Windows: chiudere le finestre cmd di start.bat)

rm -rf webui/backend/.venv      # ~150 MB
rm -rf webui/frontend/node_modules  # ~400 MB
rm -rf webui/data/timetable.db  # DATI UTENTE (salvarli prima!)
\end{lstlisting}

Per rimuovere completamente, cancellare la cartella \texttt{timetable/}.
Python, Node, Git restano installati nel sistema; rimuoverli dal
gestore di pacchetti solo se non servono ad altri progetti.

% =========== Cap 3: Architettura ===========

\chapter{Il modello a vincoli}
\label{ch:vincoli}

\epigraph{``Constraints are not the enemy of creativity. Constraints
are creativity's best friend.''}{Twyla Tharp,
\emph{The Creative Habit}}

\lettrine[lines=3]{C}{ostruire} un orario significa, dal
punto di vista formale, soddisfare un sistema di vincoli.
Alcuni di essi devono valere senza eccezione, altri esprimono
preferenze che il sistema cerca di rispettare il pi\`u
possibile pur sapendo che, in caso di conflitto, sar\`a
necessario un compromesso. Questo capitolo descrive con
ordine la tassonomia dei vincoli che $\pi$Tantum riconosce,
il modo in cui essi vengono espressi nell'interfaccia, e la
loro traduzione interna in vincoli per il solver CP-SAT.

La distinzione fra vincoli rigidi e morbidi non \`e un
dettaglio implementativo: \`e la chiave che permette al
sistema di restituire soluzioni utilizzabili. Trattare ogni
richiesta dell'utente come inderogabile produrrebbe sistemi
cronicamente infattibili, mentre trattare ogni vincolo come
puramente preferenziale violerebbe le regole contrattuali e
le esigenze didattiche di base. La gradazione in cinque
livelli che ora descriveremo \`e il risultato di un
compromesso meditato fra espressivit\`a per l'utente e
trattabilit\`a per il solver.

\section{La tassonomia dei cinque stati}

I vincoli di $\pi$Tantum sono organizzati in una tassonomia
uniforme di cinque stati, applicata sia alle matrici di
disponibilit\`a (cella per cella, per docenti, classi e
aule) sia alle preferenze di aula. I vincoli logici
disgiuntivi, espressi nel DSL del sistema, riconoscono
quattro di questi stati (l'ALLOWED non si applica a quel
contesto). La tabella~\ref{tab:cinque-stati} riassume i
cinque stati con il loro colore convenzionale, la loro
semantica e l'effetto che producono sul solver.

\begin{table}[h]
\centering
\sffamily\small
\begin{tabular}{l l p{8.5cm}}
\toprule
\textbf{Stato} & \textbf{Colore} & \textbf{Semantica} \\
\midrule
\textcolor{cAllow}{ALLOWED}  & verde chiaro & default neutro,
  applicabile soltanto alle griglie materia--aula e docente--aula;
  non contribuisce all'obiettivo \\
\textcolor{cSoft}{SOFT}      & giallo &
  penalit\`a positiva quando lo slot viene utilizzato; il valore
  \texttt{+penalty} viene sommato all'obiettivo che il solver
  minimizza \\
\textcolor{cHard}{HARD}      & rosso &
  vincolo inderogabile: il solver deve evitare lo slot oppure
  soddisfare il vincolo logico in modo assoluto \\
\textcolor{cPref}{PREFERRED} & blu &
  bonus negativo quando lo slot viene utilizzato o quando il
  vincolo logico \`e soddisfatto; il \texttt{-bonus} riduce
  l'obiettivo \\
\textcolor{cEnf}{ENFORCED}   & verde scuro &
  il solver deve forzare l'occorrenza dello slot o la
  soddisfazione del vincolo, con \emph{presence required}
  attiva \\
\bottomrule
\end{tabular}
\caption{I cinque stati della tassonomia di $\pi$Tantum, con
colori, semantica e contributo alla funzione obiettivo.}
\label{tab:cinque-stati}
\end{table}

Ognuno di questi stati ha un significato preciso che vale la
pena commentare. Lo stato ALLOWED, codificato in verde
chiaro, rappresenta la situazione di partenza: quella in cui
nessuna preferenza esplicita \`e stata espressa. \`E
applicabile soltanto a quelle griglie in cui ha senso
distinguere fra ``ammesso'' e ``vietato'', come la
compatibilit\`a fra una materia e un'aula. Per le matrici di
disponibilit\`a settimanale di docenti, classi e aule, lo
stato di partenza \`e implicitamente la disponibilit\`a, e
ALLOWED non viene mai utilizzato.

Lo stato SOFT, codificato in giallo, esprime una preferenza
negativa: l'utente preferirebbe non utilizzare quello slot,
ma se serve si pu\`o accettare. Quantitativamente, l'uso
dello slot aggiunge un costo positivo all'obiettivo, e il
solver cerca di evitarlo nei limiti del possibile. Il valore
del costo \`e configurabile, e l'utente lo regola attraverso
un campo numerico che appare nella cella stessa quando essa
viene marcata come gialla.

Lo stato HARD, codificato in rosso, esprime un divieto
assoluto. Il solver \emph{non pu\`o} utilizzare quello slot,
o equivalentemente \emph{non pu\`o} violare il vincolo
logico associato. Una soluzione che viola anche un solo
vincolo HARD non \`e una soluzione: viene scartata
immediatamente dal solver. Lo stato HARD si usa per
codificare le indisponibilit\`a contrattuali, le
incompatibilit\`a fisiche, e tutte le richieste che non
ammettono eccezioni.

Lo stato PREFERRED, codificato in blu, \`e simmetrico
rispetto al SOFT: esprime una preferenza positiva, una
configurazione che si vorrebbe vedere occorrere. L'uso
dello slot porta un bonus negativo all'obiettivo, e il
solver \`e quindi spinto a scegliere quella configurazione
quando ne ha possibilit\`a. \`E utile per esprimere
desideri come ``preferirei avere il laboratorio in due ore
consecutive'' o ``preferirei avere la palestra il
gioved\`i''.

Lo stato ENFORCED, codificato in verde scuro, esprime una
richiesta assoluta in senso opposto a HARD: il solver
\emph{deve} far accadere quella configurazione, non
soltanto deve permetterla. \`E lo strumento con cui si
codifica, per esempio, ``la riunione di consiglio si tiene
gioved\`i alla quinta ora'' (lo slot \`e
\emph{pre-allocato}). \`E un livello da usare con parsimonia,
perch\'e ENFORCED multipli e in conflitto fra loro rendono il
problema infeasibile, e il sistema non pu\`o fare altro che
segnalare l'impossibilit\`a.

\section{Le matrici di disponibilit\`a}

Le matrici di disponibilit\`a sono lo strumento principale
con cui l'utente esprime preferenze e indisponibilit\`a per
docenti, classi e aule. Ognuna di queste entit\`a possiede
una matrice di sei righe per sei colonne, che corrispondono
ai sei giorni della settimana scolastica e alle sei ore di
lezione tipiche di una giornata. Cliccando su una cella con
il mouse, l'utente la fa ciclare attraverso gli stati
disponibili nella sequenza \texttt{libero $\to$ giallo (SOFT)
$\to$ rosso (HARD) $\to$ blu (PREFERRED) $\to$ verde scuro
(ENFORCED) $\to$ libero}, in modo da poter raggiungere
qualunque stato con al massimo cinque click.

Quando una cella viene marcata come SOFT (gialla) o
PREFERRED (blu), accanto al colore appare un campo numerico
che permette di editare la penalit\`a o il bonus associati.
La logica dei segni \`e gestita automaticamente dal sistema:
i valori positivi sono ammessi soltanto sulle celle gialle e
i valori negativi soltanto su quelle blu, e il backend
applica un \emph{sign-clamp} al salvataggio per evitare
incoerenze. Trascinando il mouse mentre si tiene premuto il
tasto, si pu\`o applicare lo stesso stato a un blocco
contiguo di celle, comodit\`a che velocizza l'editing
quando un'intera giornata o un'intera fascia oraria deve
essere marcata in modo uniforme.

Una funzionalit\`a di sincronizzazione automatica governa il
rapporto fra il campo \emph{free day} del docente (il
giorno libero contrattuale) e la matrice. Quando il
\texttt{free\_day} \`e impostato, tutte e sei le ore di quel
giorno appaiono automaticamente HARD nella matrice, in modo
visivo, anche se l'utente non le ha marcate
esplicitamente. Viceversa, se l'utente edita la matrice
manualmente e tutte le ore di un giorno diventano HARD, il
backend deduce da questo che il \texttt{free\_day} \`e
quel giorno e lo aggiorna di conseguenza. La logica \`e
implementata nel modulo \texttt{webui/backend/db.py} e
garantisce che le due rappresentazioni restino sempre
coerenti.

La traduzione interna delle matrici nel modello CP-SAT
avviene nel modulo \texttt{optimization.py}, all'interno
della funzione \texttt{\_availability\_constraints}. Per
ogni entit\`a (docente, classe, aula), la matrice viene
analizzata e le sue celle si traducono in tre strutture
dati: un insieme di tuple \texttt{<entity>\_hard} che
rappresentano le \emph{forbidden assignments}, un dizionario
\texttt{<entity>\_soft} che mappa tuple a penalit\`a (con
segno gi\`a corretto), e un insieme di tuple
\texttt{<entity>\_enforced} che rappresentano le
\emph{presence required} nel solver. CP-SAT usa queste
strutture per costruire i vincoli durante la fase di
risoluzione.

\section{I vincoli logici disgiuntivi}

Accanto alle matrici di disponibilit\`a, il sistema offre
agli utenti uno strumento pi\`u espressivo per codificare
vincoli che non si lascino ridurre a singole celle: i
\emph{vincoli logici disgiuntivi}, spesso indicati con
l'acronimo DNF dall'inglese \emph{Disjunctive Normal Form}.
Si tratta di espressioni logiche che combinano slot e
predicati attraverso gli operatori AND, OR e NOT, e che si
applicano a una entit\`a (docente, classe, aula, indirizzo)
con un determinato livello di rigidit\`a (HARD, SOFT,
PREFERRED, ENFORCED).

\subsection{La grammatica delle espressioni}

La sintassi delle espressioni logiche di $\pi$Tantum \`e
stata pensata per essere compatta e leggibile. La forma BNF
semplificata \`e mostrata nel listato che segue, e descrive
gli elementi che si possono combinare.

\begin{lstlisting}[basicstyle=\ttfamily\small]
expr     := or_expr
or_expr  := and_expr ('OR' and_expr)*
and_expr := not_expr ('AND' not_expr)*
not_expr := ('NOT' | 'mai') not_expr | atom
atom     := slot | predicate | '(' expr ')'
slot     := <giorno><ora>
predicate:= <kind> ':' <name> ('@' <giorno> <ora>?)?
kind     := aula | materia | classe | gruppo | docente | prof
giorno   := lun|mar|mer|gio|ven|sab (anche lunedi|...|sabato,
            monday|...|saturday)
ora      := 1..6 (ordinale, 1=08:00) | 8..13 (assoluta)
\end{lstlisting}

La grammatica accetta anche le forme abbreviate \texttt{\&}
per AND, \texttt{|} per OR e \texttt{!} per NOT, oltre
all'alias italiano \texttt{mai} per NOT che produce
espressioni pi\`u naturali nei casi tipici. Per esempio,
\texttt{mai aula:LabFisica@mar} si legge in modo
trasparente come ``mai nel laboratorio di fisica il
marted\`i''.

\subsection{I predicate atoms}

Oltre agli slot temporali grezzi, le espressioni possono
contenere \emph{predicate atoms} che si riferiscono ad
entit\`a specifiche. La sintassi dei predicate atoms \`e
illustrata negli esempi seguenti.

\begin{lstlisting}[basicstyle=\ttfamily\small]
aula:LabFisica           -- usa l'aula LabFisica a qualunque slot
aula:LabFisica@mar       -- usa LabFisica il martedi a qualunque ora
aula:LabFisica@mar4      -- usa LabFisica martedi alla 4a ora (11:00)
materia:Religione@lun8   -- lezione di religione lunedi alle 8
classe:1A_Scientifico    -- riferimento alla classe 1A scientifico
gruppo:IRC               -- riferimento al gruppo IRC
\end{lstlisting}

I predicate atoms vengono parsati e persistiti correttamente
dal sistema. Il solver li tratta in modo coerente per le
restrizioni temporali; l'integrazione completa con la fase
di assegnazione delle aule \`e oggetto di sviluppo
incrementale, descritto nella roadmap del progetto.

\subsection{Esempi pratici}

Il modo migliore per familiarizzare con la sintassi dei
vincoli logici \`e leggere alcuni esempi reali di uso. I
seguenti esempi sono presi dalla pratica di alcune scuole
che usano $\pi$Tantum.

\begin{lstlisting}[basicstyle=\ttfamily\small]
% Per docenti
lun8 AND lun9
(lun8 AND lun9) OR (mar8 AND mar9)
gio11 OR gio12 OR gio13
NOT (mer3 AND mer4)
mai aula:LabFisica@mar
mai materia:Religione@lun8

% Per classi
mer4 AND mer5 AND mer6
mai aula:Palestra@gio
aula:LabInformatica@mar OR aula:LabInformatica@gio

% Per aule
gio4 AND gio5 AND gio6
mai materia:Religione
gruppo:IRC OR gruppo:Alternativa
\end{lstlisting}

Il primo esempio per i docenti, \texttt{lun8 AND lun9},
significa che il docente deve essere presente luned\`i alle
otto e alle nove (per esempio per una compresenza
obbligatoria). Il secondo esempio, con OR fra due
congiunzioni, ammette due alternative: o luned\`i alle
otto-nove, o marted\`i alle otto-nove. Il quarto, con NOT,
vieta la combinazione mercoled\`i terza-quarta ora. I tre
esempi con \texttt{mai} esprimono divieti su entit\`a
specifiche, leggibili in modo quasi colloquiale.

\subsection{Mapping CP-SAT}

Per ogni regola con DNF \texttt{[clause1, clause2, \ldots]},
il solver applica una traduzione che dipende dal kind del
vincolo. Una regola HARD impone che almeno una clausola sia
\emph{fully active}, ovvero che tutti i suoi atomi siano
veri. Una regola SOFT permette che nessuna clausola sia
attiva, ma in tal caso paga la penalit\`a positiva
\texttt{+soft\_penalty}. Una regola PREFERRED dota di un
bonus negativo \texttt{-soft\_penalty} la situazione in cui
almeno una clausola \`e attiva, riducendo cos\`i
l'obiettivo da minimizzare. Una regola ENFORCED si comporta
come una HARD ma con \emph{presence required}: almeno uno
degli slot della DNF deve coincidere con un'ora di lezione
attiva per l'entit\`a in questione, non basta che il vincolo
sia formalmente soddisfatto sui domini.

\section{I toggle HARD per classe}

Oltre ai vincoli espressi nelle matrici e ai vincoli logici,
il sistema mette a disposizione una serie di vincoli HARD
predefiniti, espressi come booleani nella tabella
\texttt{school\_classes}, che la maggior parte delle scuole
usa con valori standard. Sono sette: \texttt{hard\_entry\_at\_8}
(impone che la prima lezione cominci alle otto),
\texttt{hard\_exit\_after\_12} (impone che l'ultima lezione
non finisca prima delle dodici), \texttt{hard\_no\_holes}
(vieta i buchi nell'orario della classe),
\texttt{hard\_dual\_math} (impone che la matematica si
tenga in coppie di ore consecutive), \texttt{hard\_dual\_italian}
(idem per italiano), \texttt{hard\_motorie\_pairs} (idem per
educazione fisica), e \texttt{hard\_max\_6\_per\_day}
(impone il limite di sei ore al giorno).

A questi sette toggle si affianca il parametro SOFT
\texttt{soft\_minimize\_sixth\_weight}, che regola il peso
dell'obiettivo di evitare la sesta ora come ultima del
giorno. Il valore predefinito \`e tarato per esprimere una
preferenza moderata; valori pi\`u alti fanno s\`i che il
solver eviti la sesta ora con maggiore insistenza, valori
pi\`u bassi la accettano pi\`u liberamente.

\section{Il rilevamento dei conflitti}

Il tab Vincoli dell'interfaccia fornisce un bottone
denominato ``Cerca conflitti'' che, attivato, esegue una
diagnostica sul sistema di vincoli correntemente espresso e
segnala incoerenze rilevabili senza dover lanciare il
solver. Il bottone fa una chiamata HTTP a
\texttt{GET /api/monitor/conflicts}, e l'output viene mostrato
all'utente in una tabella con tre tipi di errore.

Il primo tipo, denominato \texttt{matrix\_hard\_enforced},
segnala le celle marcate sia come HARD sia come ENFORCED
contemporaneamente. Si tratta di una contraddizione
diretta, perch\'e una stessa cella non pu\`o essere allo
stesso tempo proibita e obbligatoria. Il secondo tipo,
denominato \texttt{enforced\_on\_free\_day}, segnala le
celle ENFORCED del docente che cadono nel suo giorno libero
contrattuale: anche questa \`e una contraddizione,
risolvibile o spostando l'ENFORCED o cambiando il
\texttt{free\_day}. Il terzo tipo, denominato
\texttt{logical\_unsatisfiable}, segnala i vincoli logici
HARD o ENFORCED la cui DNF non \`e soddisfacibile dato
l'insieme degli slot HARD del proprietario. \`E un caso che
richiede al coordinatore di rivedere il vincolo o di
rilassare le indisponibilit\`a.

\section{Tag delle aule}

Per esprimere caratteristiche delle aule che non si lasciano
ridurre al solo nome o tipo (``questa aula ha il
proiettore'', ``questa aula \`e adatta alla matematica
perch\'e ha la lavagna grande'', ``questa aula \`e una sala
riunioni''), il sistema offre i \emph{tag delle aule}: si
tratta di etichette libere, vere e proprie stringhe, che il
coordinatore associa a ciascuna aula dalla scheda Aule.
Ogni aula pu\`o avere quanti tag desidera, e ogni tag pu\`o
essere riutilizzato su pi\`u aule. La struttura non impone
una tassonomia gerarchica: i tag formano un semplice
insieme.

L'utilit\`a dei tag emerge soprattutto nella scrittura dei
vincoli. Invece di scrivere ``solo queste tre aule
specifiche'', enumerandole per nome (un elenco che andrebbe
aggiornato ogni volta che le aule cambiano), il coordinatore
pu\`o scrivere ``aule taggate matematica'', e il vincolo
seguir\`a automaticamente le etichettature correnti. Se
domani viene aggiunta una nuova aula taggata
\texttt{matematica}, il vincolo la includer\`a senza
ulteriori interventi.

Un esempio pratico chiarisce la logica. Supponiamo di
volere che la matematica si tenga soltanto in aule taggate
\texttt{matematica}. Una volta etichettate dalla scheda
Aule le aule che vogliamo includere, il vincolo nel DSL
generico si scrive cos\`i:

\begin{lstlisting}
forall l in lessons where l.subject == Matematica:
    "matematica" in l.classroom.tags
\end{lstlisting}

Il sistema offre anche una serie di tag automatici prodotti
dalla generazione mock. Per esempio, le aule home class
ricevono il tag dell'indirizzo della classe assegnata,
permettendo di scrivere vincoli del tipo ``le lezioni della
3A devono stare in aule taggate scientifico''. I dettagli
della generazione automatica si trovano nel modulo
\texttt{webui/backend/mock\_classrooms.py}.

\section{Le preferenze di giorno libero per i docenti}

I docenti hanno diritto, per contratto collettivo nazionale
del settore scolastico italiano, a un certo numero di giorni
liberi a settimana, tipicamente uno per il personale a
tempo pieno. $\pi$Tantum gestisce questa esigenza in modo
flessibile attraverso un meccanismo a tre slot di preferenza
ordinati. Il primo slot rappresenta il giorno libero
ideale del docente, per esempio il sabato; il secondo
rappresenta il giorno di ripiego se il primo non \`e
realizzabile, per esempio il mercoled\`i; il terzo
rappresenta l'ultima alternativa, per esempio il luned\`i.

Per ciascuna delle tre preferenze, il docente o il
coordinatore pu\`o scegliere se essa sia di tipo HARD
(assoluta, non negoziabile) o di tipo SOFT con un peso
configurabile. A questo si aggiunge il campo
\texttt{required\_free\_days\_count}, che indica quanti
giorni liberi servano in totale: il valore predefinito \`e
uno, conforme al contratto collettivo, ma scuole con
docenti a tempo parziale possono richiederne di pi\`u.

Un esempio illustra la flessibilit\`a del meccanismo. Si
consideri un docente con le preferenze [sabato HARD,
mercoled\`i SOFT con peso 50, luned\`i SOFT con peso 20] e
con \texttt{required\_free\_days\_count = 1}. Il sistema
cercher\`a sicuramente di mantenere libero il sabato,
perch\'e \`e HARD; se non riuscisse, in seconda istanza
cercher\`a di liberare il mercoled\`i (con un costo SOFT
di 50 se non ci riesce); in terza istanza cercher\`a di
liberare il luned\`i (con un costo di 20). La gradazione
permette al sistema di scegliere il compromesso pi\`u
ragionevole quando le preferenze multiple dei vari docenti
entrano in conflitto fra loro.

\section{Le preferenze di giorno libero per le classi}

Anche le classi possono avere giorni preferiti di non
attivit\`a, e la struttura del meccanismo \`e analoga a
quella dei docenti: tre slot ordinati con flag HARD o SOFT,
pi\`u un \texttt{required\_free\_days\_count}. La
motivazione \`e diversa: una classe pu\`o desiderare di
avere il sabato libero per attivit\`a di stage, oppure il
luned\`i per evitare il rientro tipicamente faticoso dopo il
fine settimana.

In aggiunta alle preferenze di giorno libero, le classi
possiedono il campo \texttt{max\_hours\_per\_day}, che
impone un tetto HARD al numero di ore in un giorno. Il
valore predefinito \`e cinque, ma alcuni istituti adottano
quattro per le classi del biennio (per ridurre il carico
giornaliero) o sei o sette per casi particolari. Riducendo
il tetto si ottengono giornate pi\`u leggere, ma se il
monte ore complessivo della classe \`e alto sar\`a
necessario estendere il calendario settimanale al sabato per
non sforare. Il sistema gestisce automaticamente questa
ridistribuzione, e segnala all'utente i casi in cui un
tetto troppo basso renderebbe il problema infeasibile.

\section{Compresenze, sostegno e potenziamento}

La scuola italiana ha tre famiglie di vincoli che il
sistema modella nativamente: le \textbf{compresenze} (il
laboratorio dove lavorano due docenti insieme), il
\textbf{sostegno} (il prof DVA che segue uno studente
specifico), e il \textbf{potenziamento} previsto dalla Legge
107 del 2015 (le ore dell'organico potenziato, senza
cattedra fissa). Sono entrati nel sistema a partire dal Task
C1 e sono enforced come constraint CP-SAT, non gestiti
post-hoc. Le violazioni sono restituite come errore al POST
con un messaggio specifico, prima che la run parta.

\subsection{Compresenze ``shared'' (lab di chimica)}

\textbf{Caso d'uso:} laboratorio di chimica del 2C, 4 ore
alla settimana, di cui 2 in compresenza con l'assistente di
laboratorio. La 2C ha le sue 4 ore di chimica, ma in 2 di
queste \`e prevista la presenza simultanea di due docenti.

\textbf{Configurazione dalla UI:}
\begin{enumerate}
\item Tab \texttt{/assignments}: la cattedra principale di
ProfChim (4 ore) \`e gi\`a presente.
\item Tab \texttt{/coteaching}: nuovo gruppo
\texttt{(2C, Chimica, n\_hours=2, kind=shared,
required=True)}.
\item Tab \texttt{/assignments}: aggiungi una nuova cattedra
ProfAss (2 ore di Chimica in 2C) e collegala al gruppo di
compresenza appena creato.
\end{enumerate}

\textbf{Cosa fa il solver:} introduce una variabile
\texttt{coday\_count[gid, d]} per ogni giorno, vincolata a
sommare a \texttt{n\_hours} sull'arco settimanale. Il
docente principale (quello con pi\`u ore, in questo caso
ProfChim) ha
\texttt{day\_count[principal, d] >= coday\_count[gid, d]}: le
sue ore di lezione coprono sempre quelle di compresenza. Il
co-docente (ProfAss) ha invece
\texttt{day\_count[codoc, d] == coday\_count[gid, d]}: le
sue 2 ore coincidono esattamente con le ore di compresenza.
A livello di slot orario, i due docenti vengono messi nella
stessa cella tramite una variabile booleana di slot
condiviso.

\subsection{Sostegno: il prof DVA che segue lo studente}

\textbf{Caso d'uso:} il prof Bianchi \`e il docente di
sostegno per uno studente DVA della 2A, e ha una cattedra
da 18 ore di sostegno. Il prof segue le lezioni della classe
(non ne aggiunge di proprie) e deve essere presente solo
quando 2A ha lezione con un altro docente.

\textbf{Configurazione:}
\begin{enumerate}
\item Tab \texttt{/teachers}: ProfBianchi esiste con la
materia ``sostegno'' nelle abilitazioni.
\item Tab \texttt{/assignments}: nuova cattedra
\texttt{(ProfBianchi, 2A, sostegno, ore=18,
is\_support=True)}.
\end{enumerate}

\textbf{Cosa fa il solver:} la cattedra di sostegno NON
contribuisce all'orario della 2A (le 18 ore non vanno a
sommarsi alle ore curricolari). Il vincolo \`e
\texttt{slot[ProfBianchi, 2A, sostegno, h] <= 2A\_busy[h]}:
il prof di sostegno \`e presente solo dove la 2A ha gi\`a
una lezione con un altro docente. La somma delle sue ore
settimanali \`e fissata a 18, distribuita in modo che il
profilo sia ragionevole.

\subsection{Potenziamento (Legge 107)}

\textbf{Caso d'uso:} la prof Verdi \`e nell'organico
potenziato e ha 18 ore settimanali di potenziamento, senza
una cattedra fissa. Le ore vengono usate per progetti,
recupero o supplenze.

\textbf{Configurazione:}
\begin{enumerate}
\item Tab \texttt{/assignments}: nuova cattedra
\texttt{(ProfVerdi, classe=NULL, materia=Potenziamento,
ore=18, is\_potenziamento=True)}. La cattedra viene
visualizzata in una sezione dedicata in fondo al tab,
\textbf{Potenziamento (Legge 107)}.
\end{enumerate}

\textbf{Cosa fa il solver:} introduce
\texttt{pot\_day\_count[ProfVerdi, d]} per ogni giorno, con
somma settimanale pari a 18. Il vincolo HARD
\texttt{pot\_day\_count[d] + ore\_di\_cattedra[d] <= 5}
garantisce che la prof non superi il tetto giornaliero. Cap
settimanale HARD: 30 ore (cinque giorni per cinque ore
ciascuno, escludendo il sabato).

\textbf{Visibilit\`a nel tab Assenze/Supplenze:} quando un
docente \`e assente, il sistema cerca un supplente fra i
disponibili. I docenti di potenziamento vengono mostrati per
primi, con un badge \textbf{POT} viola e un bordo dello
stesso colore: l'idea \`e che la loro disponibilit\`a sia
specificamente progettata per coprire questi buchi.

\section{Religione, alternativa e parallelismi intra-classe}

La classe 3B il marted\`i alla terza ora ha religione
cattolica con un docente, e contemporaneamente alternativa
con un altro docente, perch\'e gli studenti sono divisi.
Dal punto di vista dell'orario della 3B la classe \`e
``occupata'' in quello slot (gli studenti non sono
disponibili per altre lezioni), ma sono in atto due lezioni
parallele, non una.

\textbf{Configurazione (Task C2):}
\begin{enumerate}
\item Tab \texttt{/assignments}: due cattedre,
\texttt{(ProfRel, 3B, Religione, 1)} e
\texttt{(ProfAlt, 3B, Alternativa, 1)}.
\item Marca entrambe con lo stesso
\texttt{parallel\_group\_id} (id arbitrario; serve solo a
legarle insieme nel solver).
\end{enumerate}

\textbf{Cosa fa il solver:} le ore dei due docenti sono
legate slot-per-slot
(\texttt{slot[ProfRel, h] == slot[ProfAlt, h]}); a livello
di occupazione classe, le due lezioni condividono la stessa
chiave di busy, per cui la 3B risulta occupata UNA sola
volta in quello slot, anche se sono presenti due docenti.

\section{Gruppi inter-classe: seconda lingua, recupero}

Caso pi\`u complesso, introdotto dal Task C3: un gruppo che
attraversa pi\`u classi. Esempio: il gruppo di Spagnolo \`e
formato da 5 studenti della 2A e da 7 studenti della 2B,
si riunisce 3 ore alla settimana, ed \`e tenuto da
ProfSpa. Quando il gruppo si riunisce, le due classi-madre
2A e 2B sono ``occupate'' (i loro studenti sono nel gruppo,
non possono fare lezione in classe), anche se la lezione
non \`e formalmente una loro lezione.

\textbf{Configurazione:}
\begin{enumerate}
\item Tab \texttt{/students}: assicurarsi che gli studenti
del gruppo abbiano correttamente la classe-madre 2A o 2B.
\item Tab \texttt{/groups}: nuovo StudyGroup di nome
``Spagnolo 2A+2B'' con \texttt. Aggiungere i
12 studenti come membri e dichiarare le ore-materia con
\texttt{(Spagnolo, hours\_per\_week=3)}.
\item Tab \texttt{/assignments}: nuova cattedra,
selezionando il radio button \textbf{Target = gruppo
(StudyGroup)}, scegliere ``Spagnolo 2A+2B'' dal dropdown,
docente ``ProfSpa'', materia ``Spagnolo'', ore ``3''. La
cattedra appare in una sezione dedicata di colore ciano,
con badge \textbf{GRP}, separata sia dalle cattedre delle
classi che dalla sezione Potenziamento.
\end{enumerate}

\textbf{Cosa fa il solver:} per ogni ora in cui il gruppo
Spagnolo si riunisce, sia 2A che 2B vengono marcate come
occupate. L'invariante \texttt{somma di subj\_busy == pr}
del Phase B garantisce che, quando il gruppo \`e attivo in
uno slot, la classe-madre non possa avere lezioni regolari
nello stesso slot. Il vincolo per-day capacity in Phase A
\texttt{(cl\_day\_load + sum(group\_day\_count) <= 6)}
previene il caso degenere in cui curriculum e gruppo
sommerebbero pi\`u di sei ore in un solo giorno.

\subsection{Compresenza nel gruppo}

Se il gruppo Spagnolo prevede due docenti
(es. ProfSpa1 + ProfSpa2 in compresenza), si configura un
\texttt{CoteachGroup} con \texttt{group\_id} valorizzato
(XOR rispetto a \texttt{class\_id}) e si crea una seconda
cattedra con \texttt{coteach\_group\_id} sul gruppo. Il
modello CP-SAT applica la stessa logica principal/codoc che
gi\`a usa per le compresenze sulle classi.

\subsection{Sostegno nel gruppo}

Se uno studente DVA appartiene al gruppo Tedesco (3A+3B), il
prof di sostegno deve seguire lo studente \textbf{nel
gruppo} (non in classe-madre, perch\'e in quei tre slot lo
studente non \`e nella sua classe). Si configura una
cattedra
\texttt{(ProfSost, group\_id=tedesco, sostegno, ore=3,
is\_support=True)}. Il vincolo
\texttt{slot[ProfSost, h] <= group\_busy[h]} garantisce che
il prof sia presente solo nei tre slot in cui il gruppo si
riunisce, senza aggiungere ore di occupazione n\'e alla
classe-madre n\'e al gruppo.

\section{Lock nativi: blindare lezioni gi\`a sistemate}

Il pulsante a forma di lucchetto in \texttt{/assignments}
(e l'analogo nel monitor delle lezioni) marca una cattedra
come ``locked''. Quando il solver viene rilanciato per
ricalcolare l'orario, le lezioni di una cattedra lockata
sono trattate come constraint nativi: il sistema non pu\`o
spostarle da dove si trovano, deve risolvere intorno a
esse.

Vantaggio rispetto al vecchio meccanismo
``snapshot/restore'': il fail \`e immediato e parlante. Se
i lock sono incompatibili con altri vincoli (es. il giorno
libero del docente, il tetto giornaliero della classe, una
compresenza HARD), il sistema rifiuta la run con un errore
specifico, ad esempio:
\begin{quote}
\textit{Phase A INFEASIBLE: i lock attivi sono incompatibili
con i vincoli correnti. Rimuovi o adatta i lock (es.
sblocca eventi nel monitor) e ritenta. Lock attivi: 7
cattedra-giorno.}
\end{quote}

Senza lock nativi, lo stesso scenario sarebbe risultato in
una run apparentemente riuscita, con i lock silenziosamente
mossi a posteriori. I lock nativi sono propagati a tutte le
pipeline (monolitica, decomposizione temporale, spettrale,
curriculum, METIS, column generation) e a tutte le
metaeuristiche (LNS, SA, TS, ILS, ALNS, VNS, Lagrangian).

\section{Pre-flight checks: errori prima della run}

Quando l'utente lancia un'ottimizzazione, il sistema esegue
una serie di controlli sincroni \textbf{prima} di
schedulare la run. Le violazioni vengono restituite come
HTTP 400 con un elenco specifico, in italiano. Categorie
controllate:
\begin{itemize}
\item \textbf{Lock vs vincoli HARD}: lock che cadono nel
giorno libero del docente, oltre il tetto giornaliero della
classe, ecc.
\item \textbf{Compresenze}: numero di ore del principal
$\geq$ \texttt{n\_hours}; numero di ore del codoc esattamente
\texttt{n\_hours}.
\item \textbf{Sostegno}: deve avere una classe valida
(oppure un gruppo valido nel caso C3).
\item \textbf{Potenziamento}: \texttt{class\_id} deve essere
NULL; somma delle ore non deve superare 30.
\item \textbf{Gruppi (C3)}: XOR class\_id/group\_id; il
gruppo deve avere almeno uno studente; ogni studente deve
avere classe-madre; ore $>$ 0.
\end{itemize}

L'utente vede un toast di errore con la lista delle
anomalie. Pu\`o correggerle dai tab pertinenti
(\texttt{/assignments}, \texttt{/groups},
\texttt{/coteaching}) e rilanciare.


\chapter{Workflow di ottimizzazione}
\label{ch:workflow_di_ottimizzazione}

\grokfig{fig_time_budget_clessidra}{La clessidra del time budget:
sabbia a strati colorati ripartisce il tempo fra Phase~A, Phase~B
e metaeuristiche.}

\section{Schema delle fasi}

\begin{figure}[h]
\centering
\begin{tikzpicture}[node distance=0.5cm and 0.6cm,
  rect/.style={draw, rounded corners=4pt, very thick, fill=cBox!30,
               minimum width=3.4cm, minimum height=1.2cm,
               align=center, font=\small}]
  \node[rect] (a) {Phase A\\\textit{assignment}\\CP-SAT};
  \node[rect, right=of a] (b) {Phase B\\\textit{scheduling}\\+ decomp. spettrale};
  \node[rect, right=of b] (m) {Metaeuristiche\\\textit{LNS+SA+TS+ILS}};
  \node[rect, right=of m] (r) {Phase 4\\\textit{aule}\\CP-SAT};
  \draw[-{Stealth[length=3mm]}, thick] (a) -- (b);
  \draw[-{Stealth[length=3mm]}, thick] (b) -- (m);
  \draw[-{Stealth[length=3mm]}, thick] (m) -- (r);

  \node[below=0.4cm of a, font=\scriptsize] {assignments};
  \node[below=0.4cm of b, font=\scriptsize] {lessons};
  \node[below=0.4cm of m, font=\scriptsize] {lessons SOFT-ottime};
  \node[below=0.4cm of r, font=\scriptsize] {lessons.classroom};
\end{tikzpicture}
\caption{Pipeline a 4 fasi.}
\label{fig:pipeline}
\end{figure}

\section{Phase A: assegnazione (cattedre)}

\paragraph{Modulo.} \texttt{experiments/cpsat\_v2\_assignment.py}.
\paragraph{Lancio.} \texttt{POST /api/optimize/assignment} con
\texttt{\{time\_limit\_s, workers, log\}}.

\paragraph{Input.} L'insieme di \texttt{(class, subject,
hours\_per\_week)} letto da \texttt{class\_subjects} e/o
\texttt{curriculum\_subject\_hours} + le abilitazioni docente in
\texttt{teacher\_subjects} + la classe-di-concorso preference in
\texttt{subject\_group\_weights}.

\paragraph{Output.} La tabella \texttt{assignments}
(\texttt{teacher x class x subject} con \texttt{hours}). Problema di
covering: ogni \texttt{(class, subject)} riceve un docente abilitato,
rispettando max-hours per docente, classi-di-concorso e
\texttt{teacher\_compatible\_classes}.

\section{Phase B: scheduling}

\paragraph{Modulo.} \texttt{experiments/cpsat\_v2\_timetable.py} +
\texttt{experiments/decomposition\_spectral\_v2.py}.

\paragraph{Decomposizione spettrale.}
Per istanze grandi, il problema CP-SAT integrale \`e intrattabile.
La pipeline calcola un grafo di co-occupazione fra docenti, ne fa
un'embedding spettrale via Laplaciano normalizzato e usa K-means su
quell'embedding per partizionare in $k$ cluster
\emph{loosely-coupled}. Ogni cluster diventa una sub-istanza CP-SAT
indipendente; un secondo step "bridges" risolve solo i bordi; un
terzo step "ricucitura" fa full-CPSAT con tutto il problema partendo
dalla soluzione stitched.

\paragraph{Scope del solver CP-SAT (\texttt{cp\_sat\_scope}).}
Nella card 3 di \texttt{/optimize} il dropdown \emph{Scope del
solver CP-SAT} sceglie come la Phase B passa il problema al
risolutore:
\begin{itemize}
  \item \textbf{Per giorno} (\texttt{day}, default).
        Il solver risolve un giorno alla volta. La distribuzione
        settimanale delle ore di ciascuna cattedra (\emph{day count})
        viene precalcolata da Phase A e iniettata come uguaglianza
        rigida per ogni giorno. Veloce, scalabile fino a scuole grandi.
  \item \textbf{Settimana intera} (\texttt{week}).
        Il solver \texttt{MonolithicSolver(scope=None)} vede tutta
        la settimana in un solo modello CP-SAT. Pi\`u robusto sui
        vincoli che attraversano i giorni (compresenze, sostegno,
        gruppi a geometria variabile) ma significativamente pi\`u
        lento: i workers devono esplorare uno spazio di ricerca di
        sei volte la cardinalit\`a del per-giorno.
\end{itemize}

\paragraph{Modalit\`a Phase A (\texttt{phase\_a\_mode}).}
Phase A \`e il pre-step che decide quante ore di ciascuna cattedra
cadono su ciascun giorno della settimana (\emph{day count}). Il
secondo dropdown della card 3 controlla come questo output viene
combinato con il solver:
\begin{itemize}
  \item \textbf{Sempre} (\texttt{always}, default per scope = Per
        giorno). Phase A esegue e il \texttt{day\_count} viene
        imposto come uguaglianza rigida. Obbligatorio in modalit\`a
        Per giorno -- senza \texttt{day\_count} il solver per-giorno
        non sa quante ore sono "previste" in quel giorno.
  \item \textbf{Salta} (\texttt{skip}, solo con scope = Settimana
        intera). Phase A non viene eseguita; \`e la Settimana intera
        a decidere autonomamente come distribuire le ore. Massima
        libert\`a per il solver, costo: nessun warm start.
  \item \textbf{Hint morbido} (\texttt{soft\_hint}, solo con scope
        = Settimana intera). Phase A esegue come al solito ma il
        suo output viene iniettato nel solver tramite
        \texttt{model.AddHint} sui contatori per-giorno -- il solver
        parte dalla distribuzione di Phase A come prima soluzione
        ammissibile e pu\`o deviare se trova qualcosa di meglio.
        Tipicamente il compromesso preferibile per scuole complesse:
        si tiene il warm start, si toglie la rigidit\`a.
\end{itemize}

Vincolo incrociato (validato sia dal frontend sia dal backend):
\texttt{day} richiede \texttt{always}; \texttt{week} esclude
\texttt{always}. La UI disabilita le combinazioni invalide e ricorregge
automaticamente \texttt{phase\_a\_mode} quando l'utente cambia
\texttt{cp\_sat\_scope}.

\paragraph{Quando scegliere cosa.}
\begin{itemize}
  \item \emph{Scuola standard, focus sulla velocit\`a}: \texttt{day
        + always}. \`E il default per qualcosa: scala bene e produce
        soluzioni di buona qualit\`a in tempo lineare nel numero di
        classi.
  \item \emph{Scuola piccola con molte compresenze e sostegno}:
        \texttt{week + skip}. La rigidit\`a del \texttt{day\_count}
        di Phase A pu\`o rendere infeasibile un orario che invece
        esiste; lasciar decidere alla Settimana intera evita il
        problema.
  \item \emph{Scuola grande ma con vincoli incrociati notevoli}:
        \texttt{week + soft\_hint}. Il warm start velocizza la
        ricerca e la Settimana intera resta libera di adattare la
        distribuzione.
\end{itemize}

\paragraph{Branch-and-Price V1 e \texttt{cp\_sat\_scope=week}.}
Una domanda emersa durante lo sviluppo della Phase 3 era: se
Phase A non fissa pi\`u il \texttt{day\_count} come uguaglianza
rigida (\texttt{phase\_a\_mode=skip} oppure \texttt{soft\_hint}),
la generazione di colonne V1 -- master con vincoli di copertura
indicizzati su \texttt{(teacher, class, subject, day)} -- riesce a
iterare di pi\`u della singola passata
\texttt{no\_improving\_column} osservata in modalit\`a
\texttt{day}? L'intuizione era che un \texttt{dc\_value} ``allentato''
generasse duali $\lambda$ non-zero, sbloccando colonne migliorative.

La risposta empirica, misurata da
\texttt{tests/benchmarks/run\_bp\_v1\_scope\_week.py} su due scenari
sintetici (\texttt{small} = 10 docenti / 3 classi, \texttt{medium}
= 4 docenti / 3 classi con cattedre dense), \`e \emph{no}: le
cinque granularit\`a V1 (\texttt{teacher}, \texttt{teacher-day},
\texttt{teacher-class}, \texttt{teacher-class-subject},
\texttt{teacher-subject}) terminano sempre alla prima iterazione
con \texttt{bp\_terminated\_reason=no\_improving\_column} e
\texttt{bp\_columns\_added\_total=0}, sia quando il
\texttt{dc\_value} viene da Phase A sia quando viene derivato dalla
soluzione del solver settimanale.

\begin{itemize}
  \item \texttt{small}: \texttt{dc\_phase\_a} (21 chiavi non-zero)
        differisce strutturalmente da \texttt{dc\_week} (16 chiavi
        non-zero) -- 17 chiavi solo in $A$, 12 solo in $B$, 2
        chiavi in comune con valori diversi -- ma per ogni
        granularit\`a entrambe le sorgenti producono
        \texttt{iters=1, cols=0}. La soluzione \`e HARD-feasible
        in entrambi i casi; il soft score finale \`e $240$ con
        \texttt{phase\_a} e $0$ con \texttt{week}, ma la
        differenza viene dal seed pattern, non dalle iterazioni
        BP.
  \item \texttt{medium}: \texttt{dc\_phase\_a} (69 chiavi) vs
        \texttt{dc\_week} (66 chiavi), 34 vs 31 chiavi disgiunte,
        10 chiavi in comune con valori diversi. Stesso esito:
        \texttt{iters=1, cols=0} su tutte le granularit\`a, sia
        con \texttt{phase\_a} sia con \texttt{week}.
\end{itemize}

\noindent
La spiegazione \`e strutturale, non un bug della Phase 3. I seed
pattern di \texttt{column\_generation.\_seed\_patterns} sono
greedy-ottimali rispetto al \texttt{dc\_value} fornito: per ogni
docente piazzano esattamente le sue ore nei giorni indicati, in
ordine \texttt{(class, subject, day, hour)} crescente, evitando
sovrapposizioni. Il master LP V1 sceglie il pattern con peso
maggiore -- tipicamente il primo seed -- e i duali $\lambda$ del
sub-problema di pricing risultano gi\`a saturati: il pricer
ricostruisce la stessa griglia di ore e non trova \emph{reduced
cost} negativa. Cambiare la sorgente del \texttt{dc\_value}
modifica la \emph{distribuzione} delle ore (giorni e quantit\`a)
ma non cambia il fatto che ciascun docente ha un'unica ``forma''
greedy-coerente con la propria domanda. Per sbloccare iterazioni
multiple servirebbe (a) generare seed pattern \emph{strutturalmente
diversi} (perturbazioni randomiche delle ore o degli ordinamenti),
(b) introdurre vincoli di accoppiamento cross-teacher gi\`a nel
master (variante V2 con cover indicizzato su \texttt{(class, day)}
o \texttt{(curriculum, day)}), oppure (c) iniettare una pricer
heuristic con duali \emph{stabilizzati} che esplori intorno
all'ottimo locale del seed. Sono tre direzioni di lavoro distinte
e fuori dallo scope di Phase 3-5.

In pratica: \texttt{cp\_sat\_scope=week} sblocca il
\emph{solver settimanale} (la distribuzione delle ore non \`e
pi\`u prerogativa di Phase A) ma NON sblocca il loop BP V1, che
resta una single-shot ottimizzazione del seed catalog. Per orario
di buona qualit\`a su scuole grandi con BP, la leva attuale resta
la granularit\`a V2 (\texttt{class}, \texttt{class-day},
\texttt{day}, \texttt{curriculum}) o l'\texttt{iterative-diversified}
classico, non l'interazione fra scope settimanale e V1.

\section{Metaeuristiche}

\paragraph{Modulo.} \texttt{experiments/metaheuristics.py}.
Cascata di:

\begin{itemize}
  \item \textbf{LNS} -- destroy-and-repair su finestre
        \texttt{(teacher, day)} o \texttt{(class, day)}; freeze del
        60-70\% delle lezioni e CP-SAT-resolve del resto.
  \item \textbf{SA} (Simulated Annealing) -- single-swap accettati
        con Metropolis a $T$ decrescente con fattore $\alpha$.
  \item \textbf{TS} (Tabu Search) -- best-improvement local search
        con tabu list (default size 80).
  \item \textbf{ILS} (Iterated Local Search) -- alterna $n$ cicli di
        TS con \emph{kick} perturbativi (random 4-opt fra giorni).
\end{itemize}

Tutti contribuiscono al SOFT score: holes per teacher, isolated 1-
and 5-hour days, weekly distribution, dual-hour pairs, no 6th hour,
banded preferences su materie. Il SOFT include la contribuzione delle
matrici di disponibilit\`a SOFT/PREFERRED e dei logical SOFT/PREFERRED.

\section{Phase 4: assegnazione aule}

\paragraph{Modulo.} \texttt{experiments/classroom\_assignment.py}.
Input: la tabella \texttt{lessons} (g\`ia temporalmente fissata) +
le classroom rules (\texttt{kind}, \texttt{capacity},
\texttt{multi\_class}, preferenze materia$\leftrightarrow$aula,
classe$\leftrightarrow$aula, docente$\leftrightarrow$aula).
Output: \texttt{lessons.classroom\_name} popolato.

\section{Drag-and-drop con preview live}

Trascinando una lezione su un altro slot la UI fa due chiamate:

\begin{enumerate}
  \item \texttt{POST /api/schedule/move-preview} -- il backend simula
        la mossa per tutte le 36 (day, hour) e per ognuna restituisce:
        \begin{itemize}
          \item \texttt{ok} con \texttt{delta\_soft <= 0} -- mossa
                accettabile (verde se delta $<$ 0, verde chiaro
                se delta $=$ 0).
          \item \texttt{soft\_worse} -- mossa accettabile ma peggiora
                il SOFT (giallo).
          \item \texttt{hard\_violation} -- vincolo HARD violato
                (rosso).
          \item \texttt{noop} -- slot di origine.
        \end{itemize}
  \item L'utente droppa, la UI invia \texttt{PUT
        /api/schedule/move-lesson} che valida e persiste.
\end{enumerate}

\subsection{Room-follows-lesson}

Il move-preview NON considera l'aula nel calcolo HARD. Se la nuova
posizione \`e compatibile con docente / classe ma l'aula attuale
\`e g\`ia occupata o HARD-unavailable in destination, il backend:

\begin{enumerate}
  \item Sposta la lezione.
  \item Lascia l'aula libera sul nuovo slot solo se non c'\`e conflitto.
  \item Se c'\`e conflitto, rimuove l'aula dalla lezione spostata e
        risponde con \texttt{room\_cleared=true, cleared\_room=<nome>}.
  \item La UI apre un Modal che spiega cosa \`e successo e chiede di
        scegliere una nuova aula dal dropdown (verde libera / rosso
        occupata).
\end{enumerate}

\section{Slot picker matriciale (Monitor)}

Nel tab \emph{Monitor}, espandendo una cattedra si vedono le singole
lezioni; cliccando il bottone Giorno/Ora si apre un modal con una
matrice 6x6 dove ogni cella \`e colorata in base alla disponibilit\`a:

\begin{itemize}
  \item verde: free (docente e classe liberi)
  \item rosso: HARD (docente o classe g\`ia impegnati)
  \item ambra: aula occupata da un'altra lezione
  \item azzurro: slot attuale
\end{itemize}

Click su una cella libera fa un dry-run via
\texttt{PUT /api/monitor/event/\{aid\}/lesson/\{lid\}}. Se
ci sono conflitti, apre un secondo Modal con 3 opzioni:
\textbf{Annulla}, \textbf{Disassegna le lezioni in conflitto},
\textbf{Disassegna e ottimizza dopo}.

\section{Punteggio di graduatoria e criterio Anzianita}

Ciascun docente nell'anagrafica ha un campo
\textbf{graduatoria\_score} (numerico). Il valore esprime il
punteggio di graduatoria utile a discriminare i docenti quando
si tratta di assegnare cattedre o supplenze: pi\`u alto =
maggiore precedenza.

Il campo \`e usato in due punti:

\begin{enumerate}
\item \textbf{Modulo supplenze}: quando il sistema cerca un
  candidato per coprire un'assenza, ordina la rosa per
  \texttt{graduatoria\_score} decrescente.
\item \textbf{Phase A, criterio ``Anzianita''}: il DSL della
  funzione obiettivo della Phase A include un preset
  \textbf{Anzianita} che, oltre ai criteri standard di bilancio
  (capacit\`a non utilizzata, n. di indirizzi, ecc.),
  privilegia le assegnazioni che fanno coincidere docenti con
  alto \texttt{graduatoria\_score} con classi del triennio
  finale (a parit\`a di altre condizioni).
\end{enumerate}

Il preset \`e selezionabile dalla card \textbf{Phase A
(assegnazione)} del tab Workflow, dropdown ``Criterio''. Il DSL
completo della funzione obiettivo \`e visibile nel modal
``Mostra DSL attivo'' della stessa card.

% =========== Cap 6: UI guide ===========

\chapter{Guida UI}

\section{Funzionalit\`a trasversali}

\subsection{Query DSL e sort multi-livello}

Ogni pagina lista ha:
\begin{itemize}
  \item barra Query (DSL: \texttt{=, !=, <, <=, >, >=, contains,
        startswith, endswith, in [...]}, logica
        \texttt{AND / OR}, parentesi);
  \item sort multi-livello (max 4 livelli);
  \item pannello Help con la lista dei campi.
\end{itemize}

\subsection{Multi-select}

Sulle pagine Docenti / Classi / Aule la lista supporta:
\begin{itemize}
  \item click semplice: single-select (replace);
  \item Ctrl+click: toggle;
  \item Shift+click: range;
  \item bottone "Vincolo collettivo" che apre il
        \emph{BulkApplyModal}.
\end{itemize}

\subsection{Excel/CSV import}

Ogni pagina lista ha "Importa Excel/CSV" + "Template". Lo stesso
endpoint \texttt{POST /api/import/\{entity\}} gestisce le 7 entit\`a
(teachers, subjects, classes, classrooms, curricula, students,
groups). Modalit\`a: \texttt{upsert} (default), \texttt{append},
\texttt{replace}. Header alias italiani / inglesi accettati.

\subsection{Esportazione liste (xlsx / csv)}

Sopra ogni tabella, accanto a "Reset query" / "Reset sort", trovi
3 bottoni di esportazione:

\begin{itemize}
  \item \textbf{xlsx}: la \emph{vista corrente} (solo le righe
        filtrate dalla query DSL, nell'ordine determinato dal
        sort multi-livello) come file \texttt{.xlsx} con header
        colorato e auto-fit colonne.
  \item \textbf{csv}: stesso contenuto, formato CSV UTF-8 con BOM
        (compatibile con Excel italiano).
  \item \textbf{tutto}: ignora query e sort, esporta tutte le righe
        dell'entit\`a.
\end{itemize}

Filename: \texttt{<entita>\_<YYYYMMDD>\_<HHMMSS>.<ext>}. I parametri
\texttt{?format=xlsx|csv} sono accettati da tutti gli endpoint di
lista, quindi gli stessi export sono disponibili anche via curl.

\subsection{Viste salvate}

Sopra ogni tabella, accanto ai bottoni di Reset e Export, trovi un
dropdown "Viste salvate..." e i bottoni "Salva vista" / "x":

\begin{itemize}
  \item \textbf{Salva vista}: chiede un nome e memorizza la query
        DSL + il sort multi-livello come \texttt{SavedView(entity,
        name, dsl\_query, sort\_levels)}. La vista \`e globale.
  \item \textbf{Dropdown "Viste salvate..."}: applica una vista
        (sostituisce query + sort + ricarica la lista).
  \item \textbf{x}: elimina la vista correntemente selezionata.
\end{itemize}

Endpoint: \texttt{/api/saved-views?entity=<entity>} (GET / POST /
PUT \{id\} / DELETE \{id\}). Tabella: \texttt{saved\_views}.

\subsection{Tag (aule e studenti)}

Sistema many-to-many parallelo per aule e studenti.

\textbf{Tag aule}: tabella \texttt{classroom\_tags} + join
\texttt{classroom\_tag\_assignments}. UI: multi-select autocomplete
nella scheda aula, bottone "Gestisci tag" per CRUD globale,
colonna "Tag" con pill colorate (un colore stabile per nome). Il
generatore mock tagga automaticamente le aule per kind
(\texttt{lab\_fisica $\to$ [lab, fisica, scienze]}, palestra
$\to$ \texttt{[palestra, motorie]}, ecc.) e propaga il curriculum
dalla home class (\texttt{3B\_scientifico} $\to$ tag
\texttt{scientifico}).

DSL: \texttt{"matematica" in l.classroom.tags} (vincoli) o
\texttt{has\_tag(matematica)} (query lista). I tag sconosciuti
vengono creati al volo al salvataggio.

\textbf{Tag studenti}: tabella \texttt{student\_tags} + join
\texttt{student\_tag\_assignments}. Casi d'uso: \texttt{BES},
\texttt{DSA}, \texttt{debito\_matematica\_4},
\texttt{pcto\_ditta\_<x>}, \texttt{studente\_atleta}. UI analoga
a quella delle aule. Pannello "Precompila da tag" nel modal
\emph{Gruppi} che aggiunge in massa gli studenti che matchano
\texttt{any\_of} o \texttt{all\_of} su una lista di tag.

\textbf{Importante}: i tag NON sostituiscono i gruppi. I gruppi
restano l'unit\`a operativa di scheduling, i tag sono metadato
passivo per filtrare/raggruppare.

\subsection{Import bulk dedicato (\texttt{/import})}

Pagina dedicata raggiungibile dalla nav. Layout a due colonne:
configurazione (entit\`a, modalit\`a, "Scarica template") +
file da importare (drop area + selettore file). Il template xlsx
ha 3 fogli:

\begin{itemize}
  \item \texttt{Istruzioni}: tabella \texttt{Campo | Hint} con
        descrizione di ogni colonna + note generali (sinonimi
        italiani, righe vuote ignorate, etc.);
  \item \texttt{Esempi}: una riga di esempio pre-popolata;
  \item \texttt{Dati}: foglio vuoto da compilare. L'importer
        legge per default questo foglio.
\end{itemize}

Campi importabili gi\`a integrati con le feature recenti:
\texttt{tags} su \texttt{students} e \texttt{classrooms},
\texttt{graduatoria\_score} / \texttt{required\_free\_days\_count}
/ \texttt{preferred\_free\_days\_json} su \texttt{teachers}, e
\texttt{max\_hours\_per\_day} su \texttt{classes}.

\subsection{Import / Export DB (Dashboard)}

Card prominente nella prima riga della Dashboard. Funzioni:

\begin{itemize}
  \item \textbf{Esporta DB completo}: scarica un \texttt{.zip}
        contenente \texttt{database.db} (raw SQLite),
        \texttt{tables/<t>.csv} per ogni tabella e
        \texttt{metadata.json} (schema\_version + sha256).
  \item \textbf{Solo schema}: omette i dati di scheduling --
        utile come template per inizializzare un'altra scuola.
  \item \textbf{Import}: selettore file \texttt{.zip} + conferma
        esplicita. Una copia
        \texttt{timetable.db.pre\_import\_backup} viene salvata
        prima dello swap.
  \item \textbf{Snapshot}: timestampati in
        \texttt{webui/data/snapshots/}, lista con bottoni
        Ripristina / Elimina.
\end{itemize}

Endpoint: \texttt{/api/dashboard/export-db},
\texttt{/api/dashboard/import-db},
\texttt{/api/dashboard/snapshot/\{create|list|restore/<f>\}},
\texttt{DELETE /api/dashboard/snapshot/\{f\}}.

\section{Tab per tab}

\subsection*{Dashboard}
Importa profilo (5 size) con 3 checkbox per i pool
(curricula, aule, studenti); genera scuola di test; stato corrente
con 7 card-numero; \emph{RunLogPanel} streaming.

\subsection*{Docenti, Classi, Indirizzi, Studenti, Gruppi, Materie, Aule, Compresenze}
CRUD entit\`a con modal di edit. Vedere
\texttt{ui\_guide.md} per i campi specifici di ognuna.

\subsection*{Cattedre}
Layout a card: una card per classe, lista materie+docente+ore.
Bottone "cambia": apre modal con dropdown filtrato per materia,
ogni opzione mostra \texttt{Cognome Nome - assigned/max
[SFORA: +Xh]} / \texttt{[pieno]}. Bottone "Warnings cattedre":
pannello con docenti SFORA / sotto / vuoti / ok.

\subsection*{Orario}
4 viste (per classe, per docente, per aula, per slot), ognuna con
toggle matrice/lista. Drag-and-drop con preview live, dropdown
aule colorate (verde libera / rosso occupata).

\subsection*{Assenze e supplenze}
Vista settimanale 6$\times$6. Click su intestazione giorno: modal
gestione assenze. Click su cella rossa: modal con classi scoperte
(sinistra) + docenti disponibili (destra), drag-drop per assegnare
supplenti.

\subsection*{Monitor}
Lista \emph{eventi} (uno per Assignment). Espansione lezione-per-lezione
con bottone Giorno/Ora che apre slot picker matriciale (vedere
\S5.6). In cima alla pagina un \emph{segmented control} con quattro
tab per filtrare velocemente:
\textbf{Tutti}, \textbf{Incompleti} (eventi con
\texttt{n\_assigned\_hours < n\_expected\_hours}),
\textbf{Senza orario} (eventi senza alcuna lezione piazzata) e
\textbf{Lockati} (eventi con almeno una lezione marcata
\texttt{lesson.locked == true}). Per ogni riga, le azioni
disponibili sono \emph{Dissocia}, \emph{Blocca/Sblocca} e
\emph{Piazza} (eseguito singolarmente o in massa via
multi-selezione + toolbar bulk).

\subsection*{Vincoli}
Lista flat di tutti i vincoli editabili con pill colorate per livello.
Bottone "Cerca conflitti": pannello con conflitti rilevati.

\subsection*{Workflow}
Pulsanti per lanciare le 4 fasi separatamente o "Pipeline completa".
SSE log streaming.

% =========== Cap 6.bis: Tecniche di ottimizzazione ===========

\chapter{Avviare e fermare il programma}
\label{ch:launcher}
\index{Avvio}\index{Launcher}

\lettrine[lines=3]{P}{er} usare piTantum non serve digitare comandi
complicati: c'\`e uno script che avvia tutto e uno che ferma tutto.
piTantum \`e composto da due pezzi che partono insieme --- il
\emph{motore} (backend) e l'\emph{interfaccia} nel browser
(frontend) --- ma dal tuo punto di vista \`e un'unica applicazione.

\section{Avviare}

\begin{description}
\item[Windows] Fai doppio click su \texttt{webui\textbackslash
start.bat} (oppure eseguilo da terminale). Si aprono due finestre
--- lasciale aperte: sono il motore e l'interfaccia in esecuzione.
\item[Linux / macOS] Da terminale, nella cartella del progetto:
\texttt{webui/start.sh}. Parte in background; aggiungi
\texttt{-{}-foreground} se preferisci tenerlo in primo piano e
fermarlo con \texttt{Ctrl+C}.
\end{description}

Quando \`e pronto, apri il browser su \texttt{http://localhost:5173}:
\`e l'indirizzo dell'interfaccia.

\begin{avvertenzabox}
\textbf{Il primo avvio richiede qualche minuto}: la prima volta lo
script prepara da solo l'ambiente Python e le dipendenze del
frontend. Dalla seconda volta in poi l'avvio \`e immediato. Se una
delle due porte (\texttt{8000} per il motore, \texttt{5173} per
l'interfaccia) risulta gi\`a occupata, lo script te lo segnala:
probabilmente \`e rimasta aperta una sessione precedente ---
fermala prima (vedi sotto).
\end{avvertenzabox}

\section{Fermare}

\begin{description}
\item[Windows] Chiudi le due finestre aperte all'avvio.
\item[Linux / macOS] \texttt{webui/stop.sh}: ferma in modo pulito
sia il motore sia l'interfaccia.
\end{description}

\section{Cosa fanno gli script, in breve}

Entrambi gli avvii, prima di partire, fanno alcuni controlli:
verificano che Python e Node siano disponibili, creano l'ambiente
virtuale in \texttt{backend/.venv} se manca, installano le
dipendenze del frontend al primo giro, e avvisano se le porte
\texttt{8000}/\texttt{5173} sono gi\`a in uso. Su Linux/macOS
suggeriscono anche il comando giusto (\texttt{apt}, \texttt{dnf},
\texttt{pacman}, \texttt{brew}) se manca qualcosa.

\begin{biblioparvabox}
\textbf{Dettagli tecnici.} Su Windows \texttt{start.bat} usa path
relativi a \texttt{\%\textasciitilde dp0}, quindi funziona ovunque
lo si copi. Su Linux/macOS \texttt{start.sh} scrive i log in
\texttt{webui/logs/} e i PID in \texttt{webui/logs/pids};
\texttt{stop.sh} li rilegge, invia \texttt{SIGTERM} con fallback
\texttt{SIGKILL} dopo 3\,s, termina anche i processi figli
(\texttt{pkill -P}) e, come ultima rete, libera le porte
\texttt{8000}/\texttt{5173} via \texttt{lsof}.
\end{biblioparvabox}


\appendix
\chapter{Riferimenti}
Per i dettagli tecnici (architettura, modello dati, API REST,
DSL formale, strategie di ottimizzazione, diagnostica
statistica, estendere il sistema) vedere il volume separato
\textbf{Appendici tecniche} (file
\texttt{appendici\_tecniche.pdf}). Per avere entrambi i volumi
in un unico documento, vedere \texttt{manual\_completo.pdf}.

Glossario di tutti i termini tecnici: file
\texttt{docs/glossario.md} nel repository.

\end{document}
